\begin{center}
    \large

    \begin{singlespace}
        \textbf{\englishTitle{}} \\[0.5cm]
    \end{singlespace}

    \begin{singlespace}
        Student: \studentEnName\hspace{1.0cm}
        Advisor: Dr.\,\advisorEnName \\
    \end{singlespace}

    \vspace{0.5cm}
    \begin{singlespace}
        \NameofDepartmentInstituteEN\\
        National Yang Ming Chiao Tung University\\[0.2cm]
    \end{singlespace}

    \textbf{Abstract} \\[0.5cm]
\end{center}

\normalsize

Approximate nearest neighbor search (ANNS) is a core operation in vector databases, retrieval systems, and recommendation platforms.
At the hundred-million-vector scale, disk-based graph indexes such as DiskANN are widely used because the full index cannot fit entirely in memory.
However, concurrent serving creates many in-flight random SSD reads, increasing queue depth, reducing throughput, and worsening tail latency.

This thesis presents PaceANN, an acceleration layer for DiskANN-style disk-based ANNS that does not require rebuilding or relaying out the existing index.
PaceANN targets two forms of avoidable work.
First, static caches cannot adapt when hot graph nodes shift with the query distribution.
Second, fixed-budget beam search continues expanding even after easy queries have effectively converged.
To address these issues, PaceANN combines two mechanisms.
The Popularity-Adaptive Cache Manager (PAC) uses a sharded CLOCK cache to retain frequently visited nodes during serving, storing both neighbor lists and full-precision vectors to reduce traversal and re-ranking I/O.
Convergence-Aware Early Termination (CA-ET) stops each query according to its own convergence state by comparing the best unexpanded frontier candidate with exact-distance candidates already observed during search.

With 16 concurrent clients and an 8\,GB resident-memory budget, PaceANN improves throughput over a plain DiskANN baseline at comparable recall.
On SPACEV100M, it improves throughput by 1.59$\times$ at $\text{Recall@10}\approx0.98$ and reduces P99 latency by about 6\%.
On SIFT100M, it improves throughput by 1.17$\times$ at $\text{Recall@10}\approx0.99$ while keeping P99 latency comparable to the baseline.

\vspace{1cm}

\textbf{Keywords:} approximate nearest neighbor search, disk-based graph index, adaptive caching, early termination, tail latency
